% =====================================================================
\section{Casi d'uso reali}
\label{sec:casi}
% =====================================================================

\subsection{SK Telecom: orchestrazione multi-vendor su KOREN}
\label{sec:sim}

Sim et al.~\cite{sim2023} descrivono il primo dispiegamento commerciale
documentato di un orchestrator conforme alla GS~QKD~018, sul testbed della
rete di ricerca coreana KOREN (2020–2022). Due operatori QKD gestiscono due
reti con dispositivi e KMS di vendor diversi, ciascuna con il proprio
controller (QKDNC); un orchestrator (QKDNO) le coordina. I link usano BB84
con decoy state su tratte tra $39$ e $100$~km. Le due reti si toccano in
\emph{nodi di bordo} (per esempio Pangyo e Daejeon) che non contengono
moduli QKD ma due KMS, uno per operatore, interconnessi.

Quando due SAE appartengono a reti diverse, l'orchestrator calcola il
cammino end-to-end con un \emph{path computation engine}, scegliendo
\emph{a priori} il nodo di bordo da attraversare; ogni controller crea poi
il proprio segmento di service link. Per legare i segmenti gli autori
estendono il modello: introducono un tipo di applicazione
\emph{multi-segment} (\texttt{MS}), con un identificativo
\texttt{ms\_app\_id} comune ai domini e le informazioni
\texttt{server\_ms\_info} e \texttt{client\_ms\_info} sui KMS esterni, e una
RPC \texttt{app-registration-ms}. Il trasferimento delle chiavi al nodo di
bordo usa l'interfaccia KMS–KMS della GS~QKD~020~\cite{qkd020}. Il caso
mostra sia la validità dello schema 018 sia un suo limite: i tipi di
applicazione ereditati dalla 015 (\texttt{CLIENT} e \texttt{INTERNAL}) non
bastano per il multi-dominio, e servono estensioni. È un buon esempio
dell'utilità delle \texttt{identity} YANG discusse in
§\ref{sec:yang-anatomia}: un nuovo tipo si aggiunge senza riscrivere il
modello di partenza.

\subsection{Politecnico di Milano: controllo gerarchico con ONOS}

La tesi di Vicario~\cite{vicario2025} realizza un prototipo end-to-end di
controllo gerarchico per reti metropolitane: i controller di dominio sono
basati su ONOS e gestiscono i nodi tramite il modello GS~QKD~015, mentre
l'interfaccia verso l'alto è la GS~QKD~018 esposta via RESTCONF a un
orchestrator scritto in Python che aggrega la topologia di più domini. È
un'implementazione utile perché mostra l'intera pila
015~$\to$~018 in software aperto, con gli stessi moduli YANG descritti in
§\ref{sec:yang018}.

\subsection{MadQCI: la rete di Madrid}

MadQCI~\cite{madqci2024} è una rete QKD reale a Madrid con decine di moduli
di vendor diversi su oltre $130$~km, integrata con la rete ottica di
produzione. Il controllo usa uno stack SD-QKD basato sulla GS~QKD~015, che
si conferma lo standard più adottato. La GS~QKD~018 è citata come strumento
per l'orchestrazione inter-dominio e la coordinazione ottico-quantistica, ma
gli autori la descrivono come un'attività ancora in corso. È una
testimonianza onesta dello stato dell'arte: il controllo di dominio è
maturo, l'orchestrazione molto meno.
